\section{Part III --- Anatomy of the original repository}

The repository is small (about 4\,100 lines of Python). Table~\ref{tab:files} maps every file to the part of the paper it implements. Read this table twice; the modifications of Part~V are organised by the same files.

\begin{table}[h]\centering\small
\begin{tabular}{p{5.6cm}p{9.4cm}}\toprule
\textbf{File} & \textbf{What it does (paper section)} \\ \midrule
\pfile{train.py} & Entry point: loads sequences and config, builds env / model / rollout worker / data loader / evaluator; epoch loop with $\varepsilon$ and temperature schedules; checkpoints; MLL evaluation (Sec.~4.3, App.~E). \\
\pfile{src/configs/defaults.py} & All hyper-parameters (yacs/fvcore \code{CfgNode}); YAML files in \pfile{src/configs/} override them. \\
\pfile{src/env/binary_tree_env_one_step_likelihood.py} & \textbf{The MDP.} State = list of rooted subtrees; action = (pair, edge lengths); features $\rho$; reward; random rollouts; backward sampling with uniform $\PB$ (Sec.~4.1, App.~A). \\
\pfile{src/env/edges_env/continuous.py} & Random branch lengths and the action$\leftrightarrow$length conversion for the continuous model (App.~F). \\
\pfile{src/utils/evolution_model_torch.py} & JC69 transition matrices, exponential edge prior, one Felsenstein step (\code{compute\_partial\_prob}) (Sec.~3.1.1, eq.~4). \\
\pfile{src/model/transformer.py} & Transformer encoder (ViT-style blocks) over the set of lineage features (Sec.~4.3, App.~D). \\
\pfile{src/model/tree_topologies_model/one_step_model.py} & \textbf{The policy.} Embeds features, encodes the set, builds all $\binom{l}{2}$ pair features $e_i+e_j$, MLP $\to$ logits, samples a pair (App.~D). \\
\pfile{src/model/edges_model/continuous/continuous_mixture.py} & Edge MLP: Gaussian mixture over $\log b$ (one edge at the last step, two otherwise) (App.~F). \\
\pfile{src/gfn/tb_gfn_phylo.py} & \textbf{Trajectory balance.} Holds model, edge model and $\log Z$; computes the TB loss; optimiser and clipping (Sec.~3.2, eq.~3). \\
\pfile{src/gfn/rollout_worker_phylo.py} & Runs a batch of trajectories in lock-step; accumulates $\log\PF$ per step and the number of parents for $\log\PB$; $2N-3$ at the last step. \\
\pfile{src/gfn/training_data_loader.py} & Mini-batches = on-policy rollouts ($\varepsilon$-greedy) + trajectories re-sampled backward from the replay buffer of best trees (Sec.~4.3). \\
\pfile{src/gfn/gfn_evaluator.py} & Importance-weighted MLL (eq.~5) and Pearson correlation between $\log\PF^\top(x)$ and $\log\R(x)$ (Sec.~5.1). \\
\bottomrule
\end{tabular}
\caption{Files of the original repository and the paper sections they implement.}\label{tab:files}
\end{table}

\subsection{The state and the action (\file{binary\_tree\_env\_one\_step\_likelihood.py})}
A state is \code{PhylogenticTreeState(subtrees, max\_root\_idx, log\_reward)}: \code{subtrees} is a Python list of \code{PhylogeneticTree} objects (thin wrappers around \code{ete3.TreeNode}). The initial state has $N$ single-leaf trees (\code{get\_initial\_state}). The list is kept \emph{sorted by} \code{min\_seq\_idx} (the smallest leaf index in each subtree): after merging positions $i<j$, the new subtree takes position $i$ and position $j$ disappears (\code{update\_state}), which keeps the list sorted because the new subtree's minimum is the minimum of subtree $i$. This invariant is what allows the backward sampler to recompute action indices from a finished tree.

An action is a dictionary \code{\{'tree\_action': k, 'edge\_action': (b\_l, b\_r)\}}: $k$ indexes the list of all pairs $(i,j)$ of the current $l$ subtrees (\code{tree\_pairs\_dict[l]}, built once with \code{itertools.combinations}); the edge action is the pair of branch lengths of the two new edges (at the last step, where $l=2$, a single length that is split in halves, because under a reversible model only the sum matters --- the ``pulley principle'').

\subsection{The features (\code{batch\_apply\_actions}, \code{compute\_partial\_prob})}
Each subtree carries a feature tensor of shape $[m,4]$: the Felsenstein vector $L_u$ multiplied by $\prod_e P(b_e)^{1/m}$ (eq.~4 of the paper). The batch of trajectories keeps all features in one tensor \code{tree\_features} of shape $[B,l,m,4]$; \code{batch\_apply\_actions} gathers the left and right features of every trajectory, calls
\begin{lstlisting}
merged_trees_features, log_scores = self.evolution_model.compute_partial_prob(data, at_root)
\end{lstlisting}
which applies eq.~\eqref{eq:felsenstein} to the whole batch at once, multiplies by the edge priors, and returns the log-likelihood \code{log\_scores} $=\sum_i\log\sum_a L^i(a)/4$ (only meaningful at the root, but computed every step). The reward is \code{log\_rewards = (C + log\_scores) / scale} (\code{PhyloTreeReward}).

\begin{whybox}[Why the policy only sees features]
Proposition~1 of the paper: two states with the same set of features have the same optimal forward policy, so the policy need not see the tree structures at all --- only the $l$ feature vectors. This is why the model input is a \emph{set} of vectors and why a Transformer (permutation-equivariant) is used.
\end{whybox}

\subsection{The policy (\file{one\_step\_model.py})}
\begin{lstlisting}
x = self.seq_emb(batch_input)                       # [B, l, 4m] -> [B, l, 128]
x = torch.cat((summary_token, x), dim=1)            # prepend a learned [CLS]-like token
x = self.encoder(x, batch_padding_mask)             # Transformer encoder
summary_token, trees_reps = x[:, :1], x[:, 1:]
tmp = trees_reps[:, :, None, :] + trees_reps[:, None, :, :]   # e_i + e_j for all i, j
row, col = torch.triu_indices(max_nb_seq, max_nb_seq, offset=1)
x_pairs = tmp[:, row, col]                          # all distinct pairs  [B, l(l-1)/2, 128]
x_pairs = torch.cat([x_pairs, summary...], dim=2)   # concatenate the summary
logits = self.logits_head(x_pairs).squeeze(-1)      # one logit per pair
tree_actions = self.sample(ret, random_spec)        # Categorical / epsilon-greedy / temperature
log_paths_pf = log_softmax(logits)[arange(B), tree_actions]
\end{lstlisting}
The pair feature is the \emph{sum} $e_i+e_j$ (order-free), concatenated with the summary token (App.~D). \code{sample} implements $\varepsilon$-greedy exploration (\code{random\_action\_prob}) or temperature sampling (\code{T}). If \code{input\_tree\_actions} is given (replay-buffer trajectories), the sampled actions are overwritten by the given ones and only their log-probability is computed.

\subsection{Branch lengths (\file{continuous\_mixture.py})}
\code{MixtureGaussianModel} maps $[\text{summary};e_i;e_j]$ to the parameters of a Gaussian mixture over \emph{log} branch lengths ($K=5$ components; the means pass through $\tanh$ into $[-10,0]$). \code{EdgeModelContinuousMixture} uses the two-dimensional (independent) model at intermediate steps and the one-dimensional \code{root\_edge\_model} at the last step; \code{compute\_log\_path\_pf} subtracts $\log b$ (change of variables). Random exploration is not applied to branch lengths (App.~F).

\subsection{The TB loss (\file{tb\_gfn\_phylo.py})}
\begin{lstlisting}
log_pf = log_paths_pf.sum(-1)          # sum over the N-1 steps
log_pb = log_paths_pb.sum(-1)          # = -sum log |Pa(s_t)|
log_z = self.compute_log_Z(None)       # self._Z is a 256-vector; log Z = its sum (a trick for the LR)
forward_value = log_z + log_pf
backward_value = log_rewards + log_pb
loss = self.loss_fn(forward_value, backward_value)     # MSE  ==  eq. (3) averaged over the batch
\end{lstlisting}
$\log Z$ is stored as 256 numbers whose sum is $\log Z$ (so that Adam's per-parameter step is effectively 256 times larger); it is initialised from the best reward among 1\,000 random rollouts and gets its own learning rate. Temperature annealing rescales it (\code{current\_log\_z * (old\_T / new\_T)} in \file{train.py}), because $\log Z_T\approx\log Z_1/T$ when the reward is $\R^{1/T}$.

\subsection{Rollouts and the number of parents (\file{rollout\_worker\_phylo.py})}
All $B$ trajectories of a batch are advanced together (\code{while tree\_features.shape[1] > 1}), because every tree trajectory has exactly $N-1$ steps. After each step the worker records how many subtrees of the new state have children (\code{has\_children.sum(-1)}): that is $|\mathrm{Pa}(s')|$ under the ``undo one merge'' backward move. The last entry is overwritten with $2N-3$: the terminal object is unrooted, and an unrooted tree with $2N-3$ edges can be produced by $2N-3$ different last merges (one per edge on which the root could sit).

\subsection{Replay buffer and backward sampling (\file{training\_data\_loader.py}, \code{sample\_backward\_from\_tree})}
The best trees seen so far (by \code{log\_score}) are kept in a heap keyed by a signature (topology id + score). To train on one of them, the environment samples a trajectory \emph{backward}: choose uniformly one of the $2N-3$ edges as the root, then repeatedly pick uniformly a non-leaf subtree and split it (\code{sample\_backward\_prevsteps}); the forward action indices are recovered from the sorted-by-\code{min\_seq\_idx} convention; the edge actions are the stored branch lengths. The resulting action list is fed as \code{input\_actions\_set} to the rollout worker, which then only evaluates $\log\PF$ of the prescribed actions.

\subsection{Evaluation (\file{gfn\_evaluator.py})}
\code{evaluate\_marginal\_likelihood}: 1\,024 trajectories from $\PF$; \code{logsumexp(log\_scores + log\_pb - log\_pf) - log K + tree\_factor} where \code{tree\_factor} $=-\log(2N-5)!!$ is $\log P(G)$ for the uniform prior over unrooted topologies (eq.~\eqref{eq:mll}). \code{evaluate\_gfn\_quality\_pearsonr}: for 100 fixed trees, estimate $\log\PF^\top(x)$ by importance sampling over backward paths (paper eq.~S1) and correlate with $\log\R(x)$.

\subsection{Training loop (\file{train.py})}
Per epoch: compute the temperature (linear or cascading schedule), rescale $\log Z$ if it changed; build the exploration schedule for $\varepsilon$; for each of \code{STEPS\_PER\_EPOCH} batches: \code{generate\_batch} (rollouts + buffer trajectories), \code{accumulate\_loss}, \code{update\_model}; then save a checkpoint, run the evaluator, and refresh the replay buffer.
\clearpage
